\section{产品路线：Grace Blackwell、Vera Rubin 与 AI Factory}

\subsection{路线变化背后的需求信号}

Jensen 在访谈里提到，从 Grace Blackwell 到 Vera Rubin 的架构变化速度很快，原因是 workload 在变。训练之外，推理、检索、多智能体编排、企业数据接入都在改变系统瓶颈，平台必须跟着真实负载更新，而不是按固定节奏推代。

\begin{importantbox}{路线图不是时间表，而是反馈回路}
产品 roadmap 的质量，取决于能否把客户负载、前沿研究和供应链可行性实时反馈到架构决策。Jensen 用 ``listening to whispers'' 描述的就是这条反馈回路。
\end{importantbox}

\subsection{AI Factory 的单位变化}

他在后段强调，``unit of computing'' 已从单机、集群进一步变成 ``AI factory''。这意味着评估指标也要变化：不仅看训练速度，还要看 token 产能、单位能耗、运维自动化和交付周期。数据中心被当成 ``生产智能的工厂''，这是工业化视角，而不是单机视角。

\begin{knowledgebox}{AI Factory 的运营指标}
可操作的指标体系通常包括：tokens per joule、rack utilization、MTTR、job completion SLA、unit economics（每百万 token 成本）以及模型版本切换损耗。访谈虽然没有逐项展开，但 Jensen 的叙述明显指向这些运营维度。
\end{knowledgebox}

\subsection{本章小结}

NVIDIA 的产品路线已不再按 ``芯片代际'' 单轴推进，而是在 AI Factory 的多维目标下同步推进。新一代系统要回答的不是 ``快多少''，而是 ``在真实生产中稳定多快''。

